\begin{corollary}[a second determination by the spin$^c$ coboundary]\label{push:src:C.3.27}
The actual positive-loop smooth monodromy satisfies
$[\Phi]=0\in\Theta_7$. This determination uses the explicit
eight-dimensional coboundary of \ref{push:src:C.3.25} and the direct index
additivity of \ref{push:src:C.3.26}. It does not use the three-knot integer,
Wall's formula for it, the Davis generator comparison, or the
conclusion of \ref{push:src:C.3.24}.
\end{corollary}
\Uses{the hypotheses, notation and definitions stated here; \cref{push:src:C.3.25}; \cref{push:src:C.3.26}; \cref{push:src:C.3.24}; \cref{push:src:C.3.18}; \cref{push:src:C.3.22}; \cref{push:src:C.3.10}; \cref{push:src:C.3.6}; \cref{push:src:C.1.2}; \cref{push:src:C.1.1}}

\begin{proof}
Apply \ref{push:src:C.3.26} to equivariant representatives of the actual return
maps. For the annular six-loop representative $\psi$, the two
geometric characteristic calculations, with their common
orientation sign, give
\[
 p_\psi=-24\epsilon_{\rm cap},\qquad
 k_\psi=-6\epsilon_{\rm cap}.
\]
These are the actual Pontryagin and cubic pairings proved in
\ref{push:src:C.3.18}--\ref{push:src:C.3.22}. Therefore the directly defined index is
\[
 \lambda(\psi)=\frac{p_\psi-4k_\psi}{24}=0
 \quad\hbox{in }\Z.
\]

We check the relation to one positive loop before reducing
modulo $28$. Choose an invariant connection on the descended
$w$-family. Its return over six consecutive positive loops is
the sixth power of its one-loop return. A positive loop changes
$c$ by $-2$. The coherent closing translations of Paper~I, Proposition~18.5 obey
\[
 b_{h,c}+b_{k,c+h}=b_{h+k,c};
\]
in particular the six successive positive-$2$ closings from
$c_0-12$ to $c_0$ equal the positive-$12$ closing. Interpolating
between invariant connections gives equivariant isotopies.
The changes of transport in \ref{push:src:C.3.10}, its unit and real-scalar
isotopies, and the lift comparisons of \ref{push:src:C.3.6} are equivariant
as well. Changing back to the stipulated closing choices
adds only a connected $\C^*$ automorphism commuting with the
circle action. Consequently the actual lift of $\psi$ and
$\Phi^6$ define the same equivariant mapping class.

If necessary, a lift of the orbit identity on the fixed tube
is made the identity there by its smooth circle gauge; the
gauge has a real logarithm and can be extended with a collar
cutoff. This is an equivariant isotopy and does not change the
orbit map. \ref{push:src:C.3.26} therefore applies with exactly the cap used
in the calculation above, and its normalization independence
gives
\[
 6\lambda(\Phi)=\lambda(\Phi^6)=\lambda(\psi)=0.
\]
The group here is $\Z$, so $\lambda(\Phi)=0$.

Now construct the actual eight-dimensional spin$^c$ coboundary
$Z_\Phi$ of \ref{push:src:C.3.25} for a normalized one-loop representative.
Its explicit intersection form and characteristic covectors
give, by the spin$^c$ index formula alone,
\[
 \mu(\Sigma_\Phi)
   \equiv\frac{p_\Phi-4k_\Phi}{672}
   =\frac{\lambda(\Phi)}{28}=0\pmod\Z.
\]
The positive generator in \ref{push:src:C.1.2} has invariant $1/28$, and
$\Theta_7$ has order $28$. Thus this invariant detects its
elements and $[\Sigma_\Phi]=0$. The sphere-gluing comparison
of \ref{push:src:C.1.1} gives $[\Phi]=0$. It also gives
$M_\Phi\cong S^1\times S^6$ with either spin structure having
zero invariant.
\end{proof}

\begin{remark}[Independence and shared geometric inputs]
The first determination uses the integer of equivariant
orbit knots and its Wall--Davis comparison. The present
determination uses the explicit line-disc bundle and
four-handle coboundary, the spin$^c$ index formula, and an
independent integer-additivity argument from actual
seven-dimensional characteristic bordisms. Neither conversion
nor the integer-level passage from six loops to one in this
second argument invokes the first determination. The two
methods share the proved annular representative and the actual
geometric values $p_\psi=-24\epsilon_{\rm cap}$ and
$k_\psi=-6\epsilon_{\rm cap}$. Their independence is at the
boundary-invariant and one-loop arguments; they are not
calculations from disjoint geometric data. Independent
referees and exact arithmetic checks are validation, not
additional mathematical methods.

The smooth spin$^c$ coboundary here is different from the
specific analytic filling proposed in Part B. Its construction
does not supply that filling or its almost-complex Chern data.
\end{remark}
\Uses{the hypotheses, notation and definitions stated here; the argument and qualifications immediately following}


\begin{remark}[Established surgery and index machinery]
The Hopf-handle construction used in \ref{push:src:C.3.25} is established
surgery machinery. In the proof of Lemma~3.5, Case~2,
\citet{Jiang2014E} attaches the same
$S^3_{\rm Hopf}\times D^4$ four-handle to a line-disc bundle,
expresses the resulting piece as $\mathbb{CP}^2\times D^4$,
and uses its trivial normal bundle to compute the spin$^c$
invariant. For his $k=0$ case the resulting value is
$p/(24\cdot28)$; Lemma~3.6 also exhibits the characteristic
integer $(p-4k)/24$ in his oriented decomposition conventions.
\ref{push:src:C.3.25} supplies the direct calculation with general cubic
$k$ and verifies that its boundary is the actual clutched
sphere of this family. No new general surgery construction
or invariant formula is claimed.

The relation between semifree actions, fixed three-sphere
knots, and surgery to homotopy complex projective three-spaces
is classical. The primary introduction of
\citet[pp.~125--126]{MontgomeryYang1968E} describes these
correspondences and the Pontryagin offset divisible by $24$.
The second proof here derives its needed additivity and
framing independence from the explicit bordisms and the
vertical Euler identity; it does not import the knot
classification as a premise. This identifies both the
standard literature and the exact scope of the derivation.
\end{remark}
\Uses{the hypotheses, notation and definitions stated here; \cref{push:src:C.3.25}; \citet{Jiang2014E}; \citet[pp.~125--126]{MontgomeryYang1968E}; \citet{GauniyalTurchin2026Root}; \cref{push:src:C.3.17}; \citet[Theorem 2.5]{Sakai2015Root}}


\paragraph{The projection formula and its missing inputs.}
The current literature supplies another precise test for a proposed
orbit-knot computation. Theorem 4.1 of
\citet{GauniyalTurchin2026Root}, published online on
25 September 2026, expresses the Haefliger integer using both
double-point linking and invariants of the projected immersion.
The latter include a rationally normalized Smale invariant.
Consequently the double circle of \ref{push:src:C.3.17} alone is insufficient.
An application here would require an actual compressed projection
and its tangent-framing comparison. We opened the published
formula and its definitions, rather than importing the linking
term in isolation.

Likewise, \citet[Theorem 2.5]{Sakai2015Root} computes the
\emph{change} under a specified crossing change from oriented
double-preimage linking. To obtain an absolute value by that
route one must also identify the comparison embedding and its
class. Neither formula identifies the present trace with a
standard reference embedding. The characteristic-number cap
retains those geometric comparisons as explicit requirements.
Searches on 15 November 2026 for the six-sphere family,
its degeneration monodromy, and these knot formulas did not
locate a numerical determination for this particular family.
This states the search outcome.


\subsection{C.4. Change of twists and real structures}

